% !TeX root = ../../Book3.tex
\section{孤立奇点}
\label{b3:sec:0302}

奇点分类要先确认中心确实被一个完整的去心圆盘隔离。设 \(a\in\mathbb C\)、\(R>0\)，若 \(f\) 在 \(0<|z-a|<R\) 上全纯，而暂未在 \(a\) 定义，就称 \(a\) 为 \(f\) 的\term[guli-qidian]{孤立奇点}[isolated singularity]。若其他奇点不断积聚到 \(a\)，便不存在这样的去心圆盘，不能将它们全部当作中心的一个孤立奇点。

\ProseParagraphBreak

Laurent 展开把问题变成负幂系数的判别。没有负幂时，正幂级数会填入中心；只有有限个非零负幂时，乘上适当的中心距离幂就能消去分母；负幂无限多时，这两种修补都不再有效。下面分别把这三种情形转成函数值及局部因子的条件。

\subsection{可去奇点}
\label{b3:subsec:030201}

\Needspace{6\baselineskip}
\begin{definition}\label{b3:def:03-removable}
设 \(a\in\mathbb C\)、\(R>0\)，\(f:D(a,R)\setminus\{a\}\to\mathbb C\) 全纯。若存在全纯函数 \(\widetilde f:D(a,R)\to\mathbb C\)，满足 \(\widetilde f(z)=f(z)\) 对所有 \(z\ne a\) 成立，则称 \(a\) 为 \(f\) 的\term[kequ-qidian]{可去奇点}[removable singularity]。
\end{definition}

延拓只能在中心补入一个值，连续性迫使这个值等于 \(\lim_{z\to a}f(z)\)，所以若延拓存在，它必唯一。真正的判别不必事先知道这个极限，只须给出整个去心邻域的上界。

\Needspace{6\baselineskip}
\begin{theorem}[可去奇点判别]\label{b3:thm:03-removable}
设 \(a\in\mathbb C\)、\(R>0\)，\(f:D(a,R)\setminus\{a\}\to\mathbb C\) 全纯。则以下条件等价：\(a\) 是可去奇点；以 \(a\) 为中心的 Laurent 展开没有非零负幂系数；存在 \(0<r<R\)、\(M<\infty\)，使 \(|f(z)|\le M\) 对 \(0<|z-a|<r\) 成立；有限极限 \(\lim_{z\to a}f(z)\) 存在。成立时，延拓在 \(a\) 处的值为 Laurent 系数 \(c_0\)。
\end{theorem}
\begin{proof}
若 \(f\) 可去，延拓函数在 \(D(a,R)\) 内的 Taylor 展开，在去心圆盘上也是原函数的 Laurent 展开。唯一性说明所有负幂系数为零。反之，若负幂系数全为零，Laurent 展开只有正幂部分，它在整个 \(D(a,R)\) 全纯，中心值为 \(c_0\)，直接给出延拓。

若存在题设的局部上界，则对每个 \(m\ge1\) 及 \(0<\rho<r\)，\cref{b3:prop:03-laurent-coefficients}给出
\[
 |c_{-m}|\le M\rho^m.
\]
系数与半径无关，令 \(\rho\downarrow0\) 得 \(c_{-m}=0\)。这对所有 \(m\) 成立，因此有界性推出可去。可去时，延拓在一个较小闭圆盘上连续，故局部有界。最后，可去蕴含有限极限存在；有限极限又使函数在充分小的去心圆盘内有界。四个条件由此等价。
\end{proof}

例如在零点以外考虑
\[
 f(z)=\frac{e^z-1-z}{z^2}.
\]
指数级数给出 \(f(z)=1/2+z/6+z^2/24+\cdots\)，这个幂级数在零点也收敛，所以补入 \(f(0)=1/2\) 后全纯。若只沿某一条路径看极限，结论就不可靠：\(e^{-1/z^2}\) 沿实轴趋于零，但在 \(z=\mathrm{i}y\) 时等于 \(e^{1/y^2}\)，沿虚轴无界。可去判别中的上界必须同时覆盖所有充分接近中心的方向。

\ProseParagraphBreak

有时表达式本身的上界不易取得，却知道乘以 \(z-a\) 以后比常数还小。这也足以排除全部负幂。

\Needspace{6\baselineskip}
\begin{corollary}\label{b3:cor:03-removable-small-growth}
设 \(a\in\mathbb C\)、\(R>0\)，\(f:D(a,R)\setminus\{a\}\to\mathbb C\) 全纯。若
\[
 \lim_{z\to a}(z-a)f(z)=0,
\]
则 \(a\) 是 \(f\) 的可去奇点。
\end{corollary}
\begin{proof}
令 \(g(z)=(z-a)f(z)\)。有限极限使 \(g\) 可去，并使延拓在 \(a\) 处取值零。它的 Taylor 展开因而没有常数项，可写成 \(g(z)=(z-a)h(z)\)，其中 \(h\) 在整个 \(D(a,R)\) 全纯。在 \(z\ne a\) 时 \(h(z)=f(z)\)，所以 \(h\) 是所需延拓。
\end{proof}

这里极限必须等于零。函数 \(1/(z-a)\) 满足 \((z-a)f(z)=1\)，却不能补成连续函数。乘去一个分母后得到非零极限，描述的正是下一小节的一阶极点。

\subsection{极点}
\label{b3:subsec:030202}

极点可以看作全纯函数零点的倒数：去掉一个有限整数次的分母后，剩余因子在中心全纯且非零。这一非零条件确定了真正的阶数。

\Needspace{6\baselineskip}
\begin{definition}\label{b3:def:03-pole}
设 \(a\in\mathbb C\)、\(R>0\)，\(f:D(a,R)\setminus\{a\}\to\mathbb C\) 全纯。若存在整数 \(m\ge1\)、\(0<r<R\) 以及全纯函数 \(g:D(a,r)\to\mathbb C\)，满足 \(g(a)\ne0\) 且
\[
 f(z)=(z-a)^{-m}g(z)\quad(0<|z-a|<r),
\]
则称 \(a\) 为 \(f\) 的 \(m\) 阶\term[jidian]{极点}[pole]。
\end{definition}

\Needspace{6\baselineskip}
\begin{proposition}\label{b3:prop:03-pole-criterion}
设 \(a\in\mathbb C\)、\(R>0\)，\(f:D(a,R)\setminus\{a\}\to\mathbb C\) 全纯。则 \(a\) 是极点，当且仅当 \(|f(z)|\to\infty\)（\(z\to a\)）；也当且仅当 Laurent 展开只有有限个非零负幂系数且至少有一个。若最低非零指数为 \(-m\)，则极点阶数恰为 \(m\)，并有
\[
 \lim_{z\to a}(z-a)^m f(z)=c_{-m}\ne0.
\]
反之，若对某个正整数 \(m\)，上式左端存在且为非零有限值，则 \(a\) 是 \(m\) 阶极点。
\end{proposition}
\begin{proof}
若 \(f(z)=(z-a)^{-m}g(z)\)、\(g(a)\ne0\)，连续性允许缩小半径，使
\[
 \frac{|g(a)|}{2}\le|g(z)|\le\frac{3|g(a)|}{2}.
\]
因此
\[
 \frac{|g(a)|}{2|z-a|^m}\le|f(z)|
 \le\frac{3|g(a)|}{2|z-a|^m},
\]
从而 \(|f(z)|\to\infty\)。将 \(g\) 的 Taylor 展开除以 \((z-a)^m\)，最低非零指数是 \(-m\)，系数为 \(g(a)\)，其他负幂只可能有有限个，并得到所述极限。

反之，若 \(|f(z)|\to\infty\)，则在一个充分小的去心圆盘内有 \(|f(z)|>1\)，所以 \(f\) 无零点，\(h=1/f\) 全纯并趋于零。可去奇点判别使 \(h\) 延拓到中心，且 \(h(a)=0\)。它在去心圆盘上处处非零，所以不恒为零。由\cref{b3:prop:01-zero-factorization}，缩小圆盘后可写
\[
 h(z)=(z-a)^m H(z),\quad m\ge1,\quad H(a)\ne0.
\]
进一步缩小圆盘使 \(H\) 无零点，取倒数便得
\(f(z)=(z-a)^{-m}/H(z)\)，这是极点分解。

若只有有限个非零负幂系数，令 \(-m\) 为其中最低指数，则 \((z-a)^m f(z)\) 的展开只有非负幂，并在中心取非零值 \(c_{-m}\)，也给出极点分解。最后，若 \((z-a)^m f(z)\) 有非零有限极限，它由可去奇点判别延拓为中心非零的全纯函数，仍得到同一分解。最低非零指数的唯一性同时保证极点阶数唯一。
\end{proof}

这里 \(|f(z)|\to\infty\) 指对每个 \(M>0\)，都有 \(\delta>0\)，使 \(0<|z-a|<\delta\) 时 \(|f(z)|>M\)。仅仅“在任意去心邻域无界”弱得多，不能代替这一条件。本性奇点也可能在一列点上任意大，却同时在另一列点上保持有限。

\ProseParagraphBreak

商式的分母为零时，应先比较分子与分母的零点阶数。设 \(g,h\) 在 \(a\) 附近全纯且都不局部恒零，令它们在 \(a\) 的阶数分别为 \(p,q\ge0\)，其中中心非零时阶数取零。零点分解给出
\[
 g(z)=(z-a)^pG(z),\quad h(z)=(z-a)^qH(z),
 \quad G(a)H(a)\ne0.
\]
在缩小的邻域内 \(H\) 无零点，故
\[
 \frac{g(z)}{h(z)}=(z-a)^{p-q}\frac{G(z)}{H(z)}.
\]
当 \(p>q\) 时，商延拓后有 \(p-q\) 阶零点；当 \(p=q\) 时，商的可去值为 \(G(a)/H(a)\ne0\)；当 \(p<q\) 时，商有 \(q-p\) 阶极点。分母的形式阶数只有在分子不消失时才直接等于极点阶数。

\ProseParagraphBreak

例如 \(1-\cos z=z^2/2+O(z^4)\)、\(\sin z=z+O(z^3)\)，所以
\[
 \frac{1-\cos z}{z\sin z}
 =\frac{1/2+O(z^2)}{1+O(z^2)}
\]
在零点可去，补入值为 \(1/2\)。而
\[
 \frac{1-\cos z}{(\sin z)^3}
 =\frac1z\frac{1/2+O(z^2)}{1+O(z^2)}
\]
在零点只有一阶极点，尽管分母具有三阶零点。这里各个 \(O(z^2)\) 都来自完整邻域内的收敛 Taylor 级数，分母中的剩余因子在中心等于一，故可以在小圆盘内求全纯倒数。

\ProseParagraphBreak

在开集 \(U\subseteq\mathbb C\) 上，若函数除一组在 \(U\) 内没有聚点的极点外全纯，并在各极点赋值 \(\infty\)，就称它为\term[ya-chun-hanshu]{亚纯函数}[meromorphic function]。极点分解说明它局部是两个全纯函数的商，分母不局部恒零；反过来，上面的公共零点消去也说明，这样的商补入可去值后局部只有极点。例如 \(\cot z=\cos z/\sin z\) 的分母在 \(\pi\mathbb Z\) 为单零点，分子在那里不为零，故这些点都是一阶极点。

\ProseParagraphBreak

对不局部恒零的亚纯函数，可以把阶数统一约定为：零点阶数取正，极点阶数取负，非零全纯点的阶数取零。局部分解立即给出乘积阶数相加、商的阶数相减。求和时最低项可能抵消，例如
\[
 \frac1z+\left(-\frac1z+z\right)=z.
\]
两项各有极点，和却有零点。只有比较首个非零系数以后，才能判断这种抵消是否发生。

\subsection{本性奇点}
\label{b3:subsec:030203}

\Needspace{6\baselineskip}
\begin{definition}\label{b3:def:03-essential}
设 \(a\in\mathbb C\)、\(R>0\)，\(f:D(a,R)\setminus\{a\}\to\mathbb C\) 全纯。若 \(a\) 既不是可去奇点也不是极点，则称它为 \(f\) 的\term[benxing-qidian]{本性奇点}[essential singularity]。
\end{definition}

按前两项判别，本性奇点恰好对应无穷多个非零负幂系数。三类因此互斥且穷尽：主部或者全为零，或者只有有限个非零项，或者有无穷多个非零项。中心是否具有某条路径上的极限，不在这三种代数条件中。

\ProseParagraphBreak

例如
\[
 e^{1/z}=\sum_{n=0}^{\infty}\frac{z^{-n}}{n!}
 \quad(z\ne0)
\]
的负幂无限多，所以零是本性奇点。沿正实轴趋近零时函数趋于无穷，沿负实轴则趋于零。沿虚轴也没有单一极限：取
\[
 z_k=\frac1{2\pi\mathrm{i}k},
 \quad w_k=\frac1{(2k+1)\pi\mathrm{i}}
 \quad(k\ge1),
\]
两列都趋于零，却分别满足 \(e^{1/z_k}=1\)、\(e^{1/w_k}=-1\)。而指数函数没有零点，所以值零始终不被取得。一个本性奇点附近可以同时出现增长、衰减和振荡。

\ProseParagraphBreak

孤立奇点分类的前提是完整去心圆盘上已经有一个单值全纯函数。平方根在零附近的困难不满足这一前提。若存在全纯函数 \(g\) 定义在整个 \(0<|z|<R\) 上并满足 \(g(z)^2=z\)，则 \(|g(z)|=|z|^{1/2}\)，可去奇点判别使它延拓到零，且 \(g(0)=0\)。延拓不恒为零，故其零点阶数是某个正整数 \(m\)，那么 \(g^2\) 的零点阶数为 \(2m\)，不可能等于 \(z\) 的阶数一。这证明全去心圆盘上不存在这样的分支；在割去一条射线的区域选取平方根，不与此矛盾。

\subsection{Casorati–Weierstrass 定理}
\label{b3:subsec:030204}

\term[casorati-weierstrass-dingli]{Casorati--Weierstrass 定理}[Casorati--Weierstrass theorem]说明本性奇点附近的像不能避开任何开圆盘。证明的构造来自一个反问：如果函数始终离某个值有正距离，那么差的倒数就有界，可以用可去奇点判别处理。

\Needspace{6\baselineskip}
\begin{theorem}[Casorati--Weierstrass]\label{b3:thm:03-casorati-weierstrass}
设 \(a\in\mathbb C\)、\(R>0\)，\(f:D(a,R)\setminus\{a\}\to\mathbb C\) 全纯，且 \(a\) 是本性奇点。则对每个 \(0<r<R\)，集合 \(f(D(a,r)\setminus\{a\})\) 在 \(\mathbb C\) 中稠密。等价地，对任意 \(b\in\mathbb C\)、\(\varepsilon>0\) 和 \(0<r<R\)，都存在 \(z\) 满足
\[
 0<|z-a|<r,\quad |f(z)-b|<\varepsilon.
\]
\end{theorem}
\begin{proof}
若对某个 \(0<r<R\) 像集不稠密，其闭包的补集含有一个开圆盘。因此存在 \(b\in\mathbb C\)、\(\varepsilon>0\)，使
\[
 |f(z)-b|\ge\varepsilon\quad(0<|z-a|<r).
\]
令 \(g(z)=1/(f(z)-b)\)。分母无零点使 \(g\) 全纯，并且 \(|g(z)|\le1/\varepsilon\)，所以 \(g\) 可以全纯延拓到 \(a\)。

若延拓满足 \(g(a)\ne0\)，连续性保证在一个小圆盘内 \(g\) 始终非零，于是 \(b+1/g\) 在整个小圆盘全纯，并在去心部分等于 \(f\)。这使 \(a\) 可去，与假设矛盾。

若 \(g(a)=0\)，则 \(g\) 在去心圆盘上非零，因而不局部恒零。零点分解给出
\[
 g(z)=(z-a)^m h(z),\quad m\ge1,\quad h(a)\ne0.
\]
缩小圆盘使 \(h\) 无零点，再写
\[
 f(z)=b+\frac1{(z-a)^m h(z)}
 =(z-a)^{-m}\left(b(z-a)^m+\frac1{h(z)}\right).
\]
括号内是全纯函数，在中心取非零值 \(1/h(a)\)，所以 \(a\) 是 \(m\) 阶极点，也与假设矛盾。两个可能的延拓值都被排除，故像集必须稠密。
\end{proof}

定理同时允许缩小自变量邻域和目标误差。对任意给定 \(b\in\mathbb C\)，取 \(r_k=\min\{R/2,1/k\}\)、\(\varepsilon_k=1/k\)，即可选择 \(z_k\ne a\)，使 \(z_k\to a\) 且 \(f(z_k)\to b\)。它并不保证某一点真正满足 \(f(z)=b\)，例如 \(e^{1/z}\) 就遗漏零值。\chapref{b3:ch:14}的 Picard 大定理将进一步讨论精确取值，而这里的证明只给稠密性。

\ProseParagraphBreak

任意去心邻域的像既然稠密，其中的模就没有有限上界，模的下确界又必为零。这两个结论都针对整个邻域，不能强加到每一条趋近中心的路径上。稠密性给出的离散点列，也不能自动连成一条使函数值收敛的连续路径：连接两点时，函数的像可能离开指定小圆盘。点列聚值、沿路径的极限和精确取值是不同的问题。

\ProseParagraphBreak

无穷远处的孤立奇点通过换坐标归入上述分类。设 \(R_0>0\)，\(f\) 在 \(|z|>R_0\) 全纯，令
\[
 w=\frac1z,\quad F(w)=f(1/w)
 \quad(0<|w|<1/R_0).
\]
约定 \(f\) 在无穷远可去、有 \(m\) 阶极点或本性，分别是指 \(F\) 在零点具有对应类型。这里分类的是函数值；本章下一节在变换积分微分时，还会出现换元导数，不能把两种计算混为一谈。

\Needspace{6\baselineskip}
\begin{proposition}[整函数在无穷远的分类]\label{b3:prop:03-entire-at-infinity}
设 \(f:\mathbb C\to\mathbb C\) 为整函数，通过 \(F(w)=f(1/w)\) 定义它在无穷远处的奇点类型。则无穷远可去当且仅当 \(f\) 为常函数；无穷远为 \(m\) 阶极点当且仅当 \(f\) 是次数恰为 \(m\ge1\) 的多项式；无穷远为本性奇点当且仅当 \(f\) 不是多项式。
\end{proposition}
\begin{proof}
若 \(F\) 在零点可去，它在小圆盘内有界，因此 \(f\) 在某个大圆之外有界。\(f\) 在该大圆所围闭圆盘上连续，也有界，合并两项上界便得到全平面有界。由\cref{b3:thm:02-liouville}，\(f\) 为常函数。反方向直接换坐标即可。

若 \(F\) 在零点有 \(m\) 阶极点，局部分解给出小 \(|w|\) 时 \(|F(w)|\le C|w|^{-m}\)。所以大 \(|z|\) 时有 \(|f(z)|\le C|z|^m\)；再利用紧圆盘上的有界性，增加常数后得到全平面的估计
\[
 |f(z)|\le C'(1+|z|)^m.
\]
\Cref{b3:prop:02-polynomial-growth}说明 \(f\) 是次数不超过 \(m\) 的多项式。它不能是常函数，否则 \(F\) 在零点可去。若其次数为 \(d\ge1\)，最高次系数为 \(a_d\ne0\)，则
\[
 F(w)=w^{-d}\bigl(a_d+a_{d-1}w+\cdots+a_0w^d\bigr).
\]
它在零点的极点阶数恰为 \(d\)，故 \(d=m\)。这个式子也直接证明任意 \(m\) 次多项式在无穷远有 \(m\) 阶极点。最后，由孤立奇点三类的穷尽性，余下的非多项式整函数恰对应本性奇点。
\end{proof}

指数函数和正弦函数的 Taylor 展开都有无穷多个非零系数，因此它们在无穷远都是本性奇点；\(z^m\) 在无穷远则有 \(m\) 阶极点。另一方面，\(1/z\) 换成 \(w\) 后为 \(w\)，在无穷远可去且取值零，却不是常函数，因为它在有限点零处不全纯。命题中的整性负责排除这种有限奇点，不能只由大圆外的衰减推出函数恒零。
